% Source for the k = 3 "three-fold" gift exchange diagram.
% Six names arranged as a hexagon on the same 4-column, 3-row grid as
% one_component.tex / two_components.tex:
%   row1: cols 2-3, row2: cols 1 & 4, row3: cols 2-3.
% Every person has EXACTLY THREE outgoing arrows (3-regular digraph), chosen
% so the adjacency matrix is a random-looking 0/1 matrix with zero diagonal
% and every row/column summing to 3:
%   A: B C D        B: A C F        C: D E F
%   D: A E F        E: A B C        F: B D E
% The five reciprocal pairs (Alice<->Bob, Alice<->David, Bob<->Frank,
% Carlos<->Eve, David<->Frank) use bend to avoid overlap.
% CDCOLOR is the RGB triplet substituted at build time (light vs dark theme).
\documentclass{article}
\usepackage[margin=0pt]{geometry}
\usepackage{amsmath}
\usepackage{xcolor}
\usepackage{tikz-cd}
\definecolor{cdcolor}{RGB}{CDCOLOR}
\begin{document}
\thispagestyle{empty}
\begin{tikzcd}[column sep = 3em, row sep = 10 ex, arrows = {-stealth, color = cdcolor}]
  & \textcolor{cdcolor}{\text{Alice}} \arrow[from=1-2, to=1-3]
    \arrow[from=1-2, to=2-4] \arrow[from=1-2, to=3-3, bend left = 20] 
  & \textcolor{cdcolor}{\text{Bob}} \arrow[from=1-3, to=1-2, bend left = 20]
    \arrow[from=1-3, to=2-4] \arrow[from=1-3, to=2-1] & \\
  \textcolor{cdcolor}{\text{Frank}} \arrow[from=2-1, to=1-3, bend left = 20]
    \arrow[from=2-1, to=3-3, bend left = 20] \arrow[from=2-1, to=3-2]
  & & & \textcolor{cdcolor}{\text{Carlos}} \arrow[from=2-4, to=3-3]
    \arrow[from=2-4, to=3-2, bend left = 20] \arrow[from=2-4, to=2-1] \\
  & \textcolor{cdcolor}{\text{Eve}} \arrow[from=3-2, to=1-2]
    \arrow[from=3-2, to=1-3] \arrow[from=3-2, to=2-4]
  & \textcolor{cdcolor}{\text{David}} \arrow[from=3-3, to=1-2]
    \arrow[from=3-3, to=3-2] \arrow[from=3-3, to=2-1] &
\end{tikzcd}
\end{document}
